% Book 3, Year 1, chapter 12 -- program correctness.
% Provenance (never printed): sources/book3/BRIEFS.md, chapter 12.
% ledger: b1:bsearch-bug
\chapter{Program Correctness}\label{ch:b1:program-correctness}

A well-known book on programming presented a binary search, proved it correct, and tested it. The code was copied
into the standard library of a major programming language, where it ran for about nine years before it failed: on an
\omterm{def:b1:pointers-and-arrays:array}{array} of $2^{30}$ elements or more, the line \code{mid = (low + high) / 2} overflowed, and the search read outside the
\omterm{def:b1:pointers-and-arrays:array}{array}. The proof had assumed integers without limits; the machine had 32-bit ones. Proofs of programs are worth their
assumptions, and this chapter teaches how to state them: a \omterm{def:b1:program-correctness:specification}{specification} says what a program must do, a \omterm{def:b1:program-correctness:invariant}{loop
invariant} says why a loop does it, a variant says why it stops, and tests derived from the \omterm{def:b1:program-correctness:specification}{specification} check the
code against it on the cases a proof might have missed.

\begin{recall}
An \omterm{def:b1:c-modules-make-tests:assertion}{assertion} states a condition the program expects to hold at a point, and stops the program when it does not
(\cref{ch:b1:c-modules-make-tests}). A variant is a quantity in a \omterm{def:b1:induction-and-termination:well-founded}{well-founded} set that decreases at every iteration
(\cref{ch:b1:induction-and-termination}). Binary search finds a value in a sorted \omterm{def:b1:pointers-and-arrays:array}{array} by halving the part still to
search.
\end{recall}

\section{Specifications}\label{sec:b1:program-correctness:specifications}

\begin{definition}[Precondition, postcondition, specification]\label{def:b1:program-correctness:specification}
A \emph{precondition}\index{precondition} is a condition on the inputs, and on the state, that a caller must ensure
before calling a function. A \emph{postcondition}\index{postcondition} is a condition relating the result and the final
state to the inputs, which the function guarantees when it returns. The pair is the function's
\emph{specification}\index{specification}: a contract, in which the function owes nothing to a caller who broke the
precondition.
\end{definition}

\omcode{code/bachelor-1/12-program-correctness/listings/division.py}{4}{13}{Integer
division by repeated subtraction, with its contract checked}\label{lst:b1:program-correctness:division}

The docstring states the \omterm{def:b1:program-correctness:specification}{specification} in words; the \omterm{def:b1:c-modules-make-tests:assertion}{assertions} check it while the program runs. The \omterm{def:b1:program-correctness:specification}{postcondition}
$a = qb + r$ and $0 \leq r < b$ determines $q$ and $r$ completely: it is a definition of integer division, while the
loop is one way to compute it.

\omoutput{code/bachelor-1/12-program-correctness/outputs/division.txt}{Three divisions; the loop without its
\omterm{def:b1:program-correctness:specification}{precondition}, stopped by a step cap; and the \omterm{def:b1:program-correctness:specification}{precondition} refusing $b = 0$.}\label{out:b1:program-correctness:division}

\section{Partial and total correctness}\label{sec:b1:program-correctness:correctness}

\begin{definition}[Partial correctness, total correctness]\label{def:b1:program-correctness:correctness}
A program is \emph{partially correct}\index{partial correctness} with respect to a \omterm{def:b1:program-correctness:specification}{specification} if every execution
that starts in a state satisfying the \omterm{def:b1:program-correctness:specification}{precondition} and \emph{terminates} ends in a state satisfying the \omterm{def:b1:program-correctness:specification}{postcondition}.
It is \emph{totally correct}\index{total correctness} if, moreover, every such execution terminates.
\end{definition}

\omterm{def:b1:program-correctness:correctness}{Partial correctness} says nothing about the runs that never end: a loop \code{while True: pass} is \omterm{def:b1:program-correctness:correctness}{partially correct}
for every \omterm{def:b1:program-correctness:specification}{specification}. The division loop, run without its \omterm{def:b1:program-correctness:specification}{precondition}, illustrates the difference: with $b = 0$,
$r \geq b$ stays true for ever, and the loop never ends; but whenever it does end, its result satisfies the
\omterm{def:b1:program-correctness:specification}{postcondition}, so it is \omterm{def:b1:program-correctness:correctness}{partially correct} even then. \omterm{def:b1:program-correctness:correctness}{Total correctness} is \omterm{def:b1:program-correctness:correctness}{partial correctness} plus termination,
proved separately, by a variant: here $r$, which decreases by $b > 0$ at each iteration, and only because of the
\omterm{def:b1:program-correctness:specification}{precondition}.

\begin{omfigure}
\begin{tikzpicture}[s/.style={circle, fill=black, inner sep=1.2pt}, >={Stealth[length=2mm]}]
\draw[rounded corners, fill=omCode!8] (-0.4,-1.9) rectangle (1.6,1.2);
\node[font=\scriptsize\bfseries] at (0.6,1.45) {precondition holds};
\draw[rounded corners, fill=omDef!12] (6.0,-0.4) rectangle (8.4,1.2);
\node[font=\scriptsize\bfseries] at (7.2,1.45) {postcondition holds};
\node[s] (a) at (0.2,0.8) {}; \node[s] (b) at (0.2,0.0) {}; \node[s] (c) at (0.2,-0.9) {};
\node[s] (a2) at (7.0,0.6) {};
\node[s] (b2) at (7.0,-1.1) {};
\draw[->, thick] (a) .. controls (3,1.6) and (5,1.2) .. node[above, font=\scriptsize] {stops, right} (a2);
\draw[->, thick, omThm] (b) .. controls (3,-0.3) and (5,-1.2) .. node[above, sloped, font=\scriptsize, pos=0.45]
  {stops, wrong} (b2);
\draw[->, thick, dashed] (c) .. controls (2,-1.8) and (3.5,-1.4) .. (4.6,-1.7);
\node[font=\scriptsize, anchor=west] at (4.7,-1.7) {never stops};
\node[font=\scriptsize, align=left, anchor=west] at (8.7,0.4) {partial correctness:\\no run of the middle
kind\\[1mm]total correctness:\\no run of the last two kinds};
\end{tikzpicture}
\par\smallskip{\small Three kinds of runs from a state satisfying the \omterm{def:b1:program-correctness:specification}{precondition}. \omterm{def:b1:program-correctness:correctness}{Partial correctness} forbids a
run that stops in a state violating the \omterm{def:b1:program-correctness:specification}{postcondition}; \omterm{def:b1:program-correctness:correctness}{total correctness} also forbids a run that never stops.}
\end{omfigure}

\section{Loop invariants}\label{sec:b1:program-correctness:invariants}

\begin{definition}[Loop invariant]\label{def:b1:program-correctness:invariant}
A \emph{loop invariant}\index{loop invariant} is a condition on the variables of a loop that holds before the first
test of the loop condition and is preserved by each iteration: if it holds and the loop condition is true before an
iteration, it holds after it. It therefore holds at every test, and in particular when the loop exits.
\end{definition}

\begin{method}[Finding and using an invariant]\label{met:b1:program-correctness:invariant-method}
\begin{enumerate}
\item Start from the \omterm{def:b1:program-correctness:specification}{postcondition}, and weaken it into a condition that holds part of the way: \emph{drop a conjunct}
that the loop condition will supply at the exit (for division, keep $a = qb + r \wedge r \geq 0$ and drop $r < b$, which
is the negation of \code{r >= b}), or \emph{replace a constant by a variable} (for the maximum of $xs$, ``$m$ is the
maximum of $xs[0{:}i]$'' with $i$ in place of $\text{len}(xs)$).
\item \emph{Initialisation}: prove that the invariant holds after the code before the loop.
\item \emph{Preservation}: assuming the invariant and the loop condition, prove the invariant after one iteration.
\item \emph{Conclusion}: prove that the invariant and the negation of the loop condition imply the \omterm{def:b1:program-correctness:specification}{postcondition}.
\item \emph{Termination}: find a variant, and prove that it decreases at every iteration.
\end{enumerate}
\end{method}

\omcode{code/bachelor-1/12-program-correctness/listings/maximum.py}{4}{13}{The maximum, with its invariant as an
\omterm{def:b1:c-modules-make-tests:assertion}{assertion}}\label{lst:b1:program-correctness:maximum}

\omoutput{code/bachelor-1/12-program-correctness/outputs/maximum.txt}{Three maxima, every invariant check
passed.}\label{out:b1:program-correctness:maximum}

The \omterm{def:b1:c-modules-make-tests:assertion}{assertion} turns the invariant into a check that runs at every iteration, which costs time (here it makes the loop
quadratic) and is for development and tests only. The proof itself is \cref{exo:b1:program-correctness:8}. For the
division, the obligations are short. Initialisation: $q = 0$ and $r = a$ give $a = 0 \cdot b + a$ and $r = a \geq 0$.
Preservation: if $a = qb + r$ and $r \geq b$, then $a = (q + 1)b + (r - b)$ and $r - b \geq 0$. Conclusion: with $r < b$
at the exit, the \omterm{def:b1:program-correctness:specification}{postcondition} holds. Termination: $r$ is a natural number that decreases by $b \geq 1$.

\section{Proofs for array algorithms}\label{sec:b1:program-correctness:arrays}

\subsection*{Partition}

Partitioning around a pivot rearranges a slice so that the elements smaller than the pivot come first, then the
pivot, then the others; quicksort is built on it in \cref{ch:b1:sorting}. Lomuto's version scans the slice
once, with two indices.

\omcode{code/bachelor-1/12-program-correctness/listings/partition.py}{4}{17}{Lomuto's
partition}\label{lst:b1:program-correctness:partition}

\omoutput{code/bachelor-1/12-program-correctness/outputs/partition.txt}{The three regions before each step: smaller,
not smaller, not yet seen.}\label{out:b1:program-correctness:partition}

\begin{omfigure}
\begin{tikzpicture}[x=7mm, y=7mm]
\fill[omDef!20] (0,0) rectangle (3,1); \fill[omCode!12] (3,0) rectangle (6,1); \fill[black!4] (6,0) rectangle (10,1);
\fill[omThm!20] (10,0) rectangle (11,1);
\draw (0,0) rectangle (11,1); \foreach \i in {1,...,10} {\draw (\i,0) -- (\i,1);}
\node[font=\small] at (1.5,1.4) {$< p$}; \node[font=\small] at (4.5,1.4) {$\geq p$};
\node[font=\small] at (8,1.4) {not yet seen}; \node[font=\small] at (10.5,1.4) {$p$};
\foreach \x/\t in {0/lo, 3/k, 6/j, 10/hi-1, 11/hi} {\draw[thick] (\x,-0.15) -- (\x,0); \node[font=\scriptsize\ttfamily,
  below] at (\x,-0.15) {\t};}
\node[font=\scriptsize, align=left, anchor=west] at (11.6,0.5) {if $a[j] < p$: swap $a[k]$ and $a[j]$,\\then $k$
and $j$ move right;\\otherwise only $j$ moves right};
\end{tikzpicture}
\par\smallskip{\small The invariant of Lomuto's partition: $a[lo{:}k] < p$ and $a[k{:}j] \geq p$, the rest unseen. A
step looks at $a[j]$ and extends one of the two regions; when $j$ reaches $hi - 1$, the pivot is swapped into place
$k$, between them.}
\end{omfigure}

\begin{proposition}[Lomuto's partition is correct]\label{prop:b1:program-correctness:partition}
If $lo < hi$, \code{partition(a, lo, hi)} terminates, permutes $a[lo{:}hi]$, and returns $k$ with $a[lo{:}k] < a[k]$ and
$a[k + 1{:}hi] \geq a[k]$, where $a[k]$ is the original $a[hi - 1]$.
\end{proposition}

\begin{proof}
The loop runs $j$ from $lo$ to $hi - 2$, so it terminates, and the only changes to $a$ are swaps: a permutation.
Invariant: $a[lo{:}k] < p$ and $a[k{:}j] \geq p$. At the start, $k = j = lo$, both regions are empty. If $a[j] \geq p$,
only $j$ moves: the second region gains $a[j]$. If $a[j] < p$, swapping $a[k]$ and $a[j]$ puts $a[j]$ at the end of the
first region and moves the first element of the second region (if any) to its end; then $k$ and $j$ move. At the exit,
$j = hi - 1$: the slice is $< p$ on $[lo, k)$, $\geq p$ on $[k, hi - 1)$, and $p$ at $hi - 1$. The final swap puts $p$ at
$k$ and the element $a[k] \geq p$ (if $k < hi - 1$) at $hi - 1$, inside the second region.
\end{proof}

\subsection*{Binary search on a half-open interval}

The half-open interval $[lo, hi)$, the integers $i$ with $lo \leq i < hi$, is the convenient way to describe a slice:
its length is $hi - lo$, it is empty when $lo = hi$, and $[lo, mid)$ and $[mid, hi)$ split it without overlap. Binary
search keeps such an interval of candidates.

\omcode{code/bachelor-1/12-program-correctness/listings/bsearch.py}{4}{15}{The leftmost position of $x$, on
$[lo, hi)$}\label{lst:b1:program-correctness:bsearch}

\omoutput{code/bachelor-1/12-program-correctness/outputs/bsearch.txt}{Present, with duplicates, at both ends, absent
inside and outside, and the empty list.}\label{out:b1:program-correctness:bsearch}

\begin{proposition}[Binary search is correct]\label{prop:b1:program-correctness:bsearch}
If $a$ is sorted, \code{leftmost(a, x)} terminates and returns the least $i$ with $a[i] \geq x$, or $\text{len}(a)$ if
there is none.
\end{proposition}

\begin{proof}
Invariant: $0 \leq lo \leq hi \leq n$, every element of $a[{:}lo]$ is $< x$, and every element of $a[hi{:}]$ is $\geq
x$. Initially $lo = 0$ and $hi = n$, and both slices are empty. In an iteration $lo < hi$, so $lo \leq mid < hi$: $mid$
is a valid index, and the interval shrinks whichever branch runs. If $a[mid] < x$, then, $a$ being sorted, all of
$a[{:}mid + 1]$ is $< x$, and $lo = mid + 1$ keeps the invariant; otherwise all of $a[mid{:}]$ is $\geq x$, and $hi = mid$
keeps it. Variant: $hi - lo$, which decreases since $mid + 1 > lo$ and $mid < hi$. At the exit $lo = hi$: everything
before $lo$ is $< x$ and everything from $lo$ on is $\geq x$, so $lo$ is the least index with $a[lo] \geq x$, or $n$.
\end{proof}

\begin{omfigure}
\begin{tikzpicture}[x=6mm, y=6mm]
\fill[omDef!20] (0,0) rectangle (4,1); \fill[black!4] (4,0) rectangle (9,1); \fill[omCode!12] (9,0) rectangle (13,1);
\draw (0,0) rectangle (13,1); \foreach \i in {1,...,12} {\draw (\i,0) -- (\i,1);}
\node[font=\small] at (2,1.4) {$< x$}; \node[font=\small] at (6.5,1.4) {candidates}; \node[font=\small] at (11,1.4)
{$\geq x$};
\foreach \x/\t in {0/0, 4/lo, 6/mid, 9/hi, 13/n} {\draw[thick] (\x,-0.15) -- (\x,0); \node[font=\scriptsize\ttfamily,
  below] at (\x,-0.15) {\t};}
\node[font=\scriptsize, align=left, anchor=north west] at (0,-1.0) {boundary cases: empty list ($n = 0$: no
iteration, result 0); $x$ below every element (result 0);\\$x$ above every element (result $n$); $x$ present several
times (its first position);\\$x$ absent between two elements (where to insert it, keeping the list sorted)};
\end{tikzpicture}
\par\smallskip{\small The invariant of binary search on $[lo, hi)$: everything left of $lo$ is smaller than $x$,
everything from $hi$ on is at least $x$, and the answer lies in $[lo, hi]$. Each step moves one end past $mid$.}
\end{omfigure}

At most $\lfloor \log_2 n \rfloor + 1$ iterations run (\cref{exo:b1:program-correctness:9}). The midpoint is written
$lo + (hi - lo)/2$ rather than $(lo + hi)/2$: in Python both are correct, but in C the sum of two large \code{int}
indices overflows (\cref{ch:b1:c-types-and-control}), which is exactly the failure of the opening story. The
difference between the two versions is invisible to the proof unless the proof states that the variables are
mathematical integers.

\subsection*{The three-colour flag}

A list holds only the values 0, 1 and 2 (three colours); sort it in one pass, with swaps only. This is the
\emph{three-colour flag problem}\index{three-colour flag problem}, posed by Edsger Dijkstra. Its invariant has four
regions.

\omcode{code/bachelor-1/12-program-correctness/listings/flag.py}{4}{19}{The three-colour flag in one
pass}\label{lst:b1:program-correctness:flag}

\omoutput{code/bachelor-1/12-program-correctness/outputs/flag.txt}{Nine values sorted, with six
swaps.}\label{out:b1:program-correctness:flag}

\begin{omfigure}
\begin{tikzpicture}[x=6mm, y=6mm]
\fill[omDef!20] (0,0) rectangle (3,1); \fill[omCode!12] (3,0) rectangle (6,1); \fill[black!4] (6,0) rectangle (10,1);
\fill[omThm!20] (10,0) rectangle (13,1);
\draw (0,0) rectangle (13,1); \foreach \i in {1,...,12} {\draw (\i,0) -- (\i,1);}
\node[font=\small] at (1.5,1.4) {0}; \node[font=\small] at (4.5,1.4) {1}; \node[font=\small] at (8,1.4) {unknown};
\node[font=\small] at (11.5,1.4) {2};
\foreach \x/\t in {0/0, 3/low, 6/mid, 10/high, 13/n} {\draw[thick] (\x,-0.15) -- (\x,0); \node[font=\scriptsize\ttfamily,
  below] at (\x,-0.15) {\t};}
\node[font=\scriptsize, align=left, anchor=west] at (13.6,0.5) {$a[mid] = 0$: swap with $a[low]$, low and mid
advance\\$a[mid] = 1$: mid advances\\$a[mid] = 2$: high retreats, swap with $a[high]$};
\end{tikzpicture}
\par\smallskip{\small The four regions of the three-colour flag and the three cases of a step. Each step shrinks the
unknown region by one, so $high - mid$ is a variant; when it is empty, the list is sorted.}
\end{omfigure}

In the case 2, \code{mid} does not advance: the value swapped in from \code{high} has not been examined yet. In the
case 0, it may advance, because the value swapped out from \code{low} came from the region of 1s (or is the same
element, when $low = mid$).

\section{Testing against the specification}\label{sec:b1:program-correctness:testing}

A proof covers every input but rests on assumptions and can contain errors; tests run the real code on the real
machine but cover only the inputs they try. The \omterm{def:b1:program-correctness:specification}{specification} connects them: it says what any correct output is, so
it can judge outputs automatically.

\begin{definition}[Test oracle, property-based test]\label{def:b1:program-correctness:oracle}
A \emph{test oracle}\index{test oracle} is a procedure that decides whether an output is correct for an input: a
check of the \omterm{def:b1:program-correctness:specification}{postcondition}, or a comparison with a simpler program known to be correct (a brute-force or library
version). A \emph{property-based test}\index{property-based test} states a property that must hold for all inputs of
a kind, generates many inputs at random, checks the property on each, and on a failure reports a smallest failing
input it can find.
\end{definition}

\omcode{code/bachelor-1/12-program-correctness/tests/test_bsearch.py}{18}{26}{Two oracles for binary search: the
library and the \omterm{def:b1:program-correctness:specification}{postcondition}}\label{lst:b1:program-correctness:test}

The library \code{hypothesis} generates the sorted lists and the values; the standard module \code{bisect} serves as
the first oracle, and the \omterm{def:b1:program-correctness:specification}{postcondition}, checked element by element, as the second. A test is only as good as its
inputs: these lists hold small values, so that duplicates and absent values are frequent, which is where binary
searches fail. No test, however many inputs it tries, would have found the overflow of the opening story on a machine
with too little memory for an \omterm{def:b1:pointers-and-arrays:array}{array} of $2^{30}$ elements: tests can show the presence of errors, and proofs their
absence under stated assumptions.

\begin{safety}
An \omterm{def:b1:c-modules-make-tests:assertion}{assertion} that checks an invariant can be expensive and can be disabled (\cref{ch:b1:c-modules-make-tests}); it
must never have a \omterm{def:b1:c-types-and-control:side-effect}{side effect} the program relies on. Inputs from users or files are checked by code that reports an
error, not by \omterm{def:b1:c-modules-make-tests:assertion}{assertions}: a violated \omterm{def:b1:program-correctness:specification}{precondition} at a public boundary is the caller's error, and the program must
say so.
\end{safety}

\section{Exercises}\label{sec:b1:program-correctness:exercises}

\begin{exercise}[$\star$]\label{exo:b1:program-correctness:1}
\texttt{def f(a, b): return a - a \% b} is called with integers. Which of these \omterm{def:b1:program-correctness:specification}{specifications} does it meet? (a) pre
$a \geq 0, b > 0$; post $b$ divides the result. (b) pre $a \geq 0, b > 0$; post $a - b < f(a, b) \leq a$. (c) pre $b
\neq 0$; post the result is at most $a$. (d) pre $b > 0$; post the result is a multiple of $b$, at most $a$, and the
largest such.
\end{exercise}

\begin{exercise}[$\star$]\label{exo:b1:program-correctness:2}
Give the values of $i$ and $m$ at each test of the loop of \code{maximum([3, 9, 2, 9, 4])}, and check the invariant at
each.
\end{exercise}

\begin{exercise}[$\star$]\label{exo:b1:program-correctness:3}
Does each loop terminate for every natural number $n$? Give a variant or a run that does not stop.
(a) \code{while n != 0: n = n - 2}. (b) \code{while n > 1: n = n // 2}. (c) \texttt{while n > 0: n = n - 1 if n \% 2 else
n + 1}. (d) \texttt{while n \% 7 != 0: n = n + 3}.
\end{exercise}

\begin{exercise}[$\star$]\label{exo:b1:program-correctness:4}
Predict what the program prints.
\omcode{code/bachelor-1/12-program-correctness/listings/exo_dups.py}{3}{6}{Duplicates and the leftmost position}
\end{exercise}

\begin{exercise}[$\star\star$]\label{exo:b1:program-correctness:5}
This search uses a closed interval $[lo, hi]$ but updates it as a half-open one. Find the inputs on which it never
stops, explain why with the variant, and correct it with its invariant.
\omcode{code/bachelor-1/12-program-correctness/listings/bsearch_bug.py}{4}{18}{A search that can loop}
\omoutput[2-3]{code/bachelor-1/12-program-correctness/outputs/bsearch_bug.txt}{Two answers, then a step cap.}
\end{exercise}

\begin{exercise}[$\star\star$]\label{exo:b1:program-correctness:6}
A student writes the partition with a copy instead of a swap. Which part of the proof of
\cref{prop:b1:program-correctness:partition} fails, and what happens to the list?
\omcode{code/bachelor-1/12-program-correctness/listings/exo_partition_bug.py}{4}{12}{A partition with a copy}
\omoutput[2-3]{code/bachelor-1/12-program-correctness/outputs/exo_partition_bug.txt}{The result.}
\end{exercise}

\begin{exercise}[$\star\star$]\label{exo:b1:program-correctness:7}
Write \code{rightmost(a, x)}, which returns the least $i$ with $a[i] > x$ in a sorted list, on a half-open interval,
with its invariant as a comment. Deduce the number of occurrences of $x$, and test both against \code{bisect}.
\end{exercise}

\begin{exercise}[$\star\star$]\label{exo:b1:program-correctness:8}
Prove that \code{maximum} is \omterm{def:b1:program-correctness:correctness}{totally correct}, with the invariant of \cref{lst:b1:program-correctness:maximum} and the
variant $\text{len}(xs) - i$.
\end{exercise}

\begin{exercise}[$\star\star$]\label{exo:b1:program-correctness:9}
Prove that binary search on $n \geq 1$ elements makes at most $\lfloor \log_2 n \rfloor + 1$ iterations, and that the
bound is reached.
\end{exercise}

\begin{exercise}[$\star\star\star$]\label{exo:b1:program-correctness:10}
The flag of \cref{lst:b1:program-correctness:flag} sorts in one pass, but not with the fewest swaps. Find a list on which
it makes more swaps than necessary. Then write \code{min_swaps(a)}, the least number of swaps that sort a list of 0, 1
and 2, and check it against a search over all orders for lists of up to six values.
\end{exercise}

\begin{exercise}[$\star\star\star$]\label{exo:b1:program-correctness:11}
Find what this function computes, state its \omterm{def:b1:program-correctness:specification}{specification}, and prove it \omterm{def:b1:program-correctness:correctness}{totally correct} with an invariant relating
\code{acc}, \code{base} and \code{e} to $x$ and $n$. How many iterations does it make?
\omcode{code/bachelor-1/12-program-correctness/listings/exo_mystery.py}{4}{12}{A loop to understand}
\end{exercise}

\begin{exercise}[$\star\star\star$]\label{exo:b1:program-correctness:12}
A student proves that the loop \code{s = 0; i = 0; while i < n: i = i + 1; s = s + i} computes $s = n(n + 1)/2$ with the
invariant ``$s = i(i + 1)/2$ and $i \leq n$'', and proves the conclusion as follows: ``at the exit $i \leq n$ and $i \geq n$,
so $i = n$''. The same student proves that \code{while i != n: i = i + 2} (from $i = 0$) ends with $i = n$, with the
invariant ``$i$ is even''. Which proof is right, which is wrong, and where exactly?
\end{exercise}

\section{Problem: The square root that never lies}\label{sec:b1:program-correctness:problem}

\begin{problem}[{Weekend problem --- The square root that never lies}]\label{pb:b1:program-correctness:1}
The integer square root of a natural number $n$ is the $r$ with $r^2 \leq n < (r + 1)^2$. Two programs compute it:
bisection on the half-open interval $[0, n + 1)$, and Newton's method on integers, started at $n$. Both are compared
with the library's \code{math.isqrt} on twelve values up to $2^{64} - 1$, and by a \omterm{def:b1:program-correctness:oracle}{property-based test} on random 64-bit
values.
\omcode{code/bachelor-1/12-program-correctness/listings/isqrt.py}{7}{28}{Two integer square
roots}\label{lst:b1:program-correctness:isqrt}
\omoutput{code/bachelor-1/12-program-correctness/outputs/isqrt.txt}{Equal roots, and the iterations of each
method.}\label{out:b1:program-correctness:isqrt}

\medskip\noindent\textbf{Part I --- The specification.}
\begin{enumerate}
\item Show that for each $n \geq 0$ exactly one natural number $r$ satisfies $r^2 \leq n < (r + 1)^2$.
\item Show that the \omterm{def:b1:program-correctness:specification}{specification} can be checked with two multiplications, without computing any square root.
\item Why is $\lfloor \sqrt{n} \rfloor$, computed in floating point, not a correct implementation for 64-bit $n$? Give
the order of magnitude of $n$ beyond which \code{float} cannot even represent every integer.
\item Give the \omterm{def:b1:program-correctness:specification}{specification} of \code{isqrt_bisect} and \code{isqrt_newton} as \omterm{def:b1:program-correctness:specification}{precondition} and \omterm{def:b1:program-correctness:specification}{postcondition},
counting the returned number of steps.
\item Write a \omterm{def:b1:program-correctness:oracle}{test oracle} for the integer square root that uses only the \omterm{def:b1:program-correctness:specification}{specification}.
\item Why do the programs return the step counts with the result?
\end{enumerate}

\medskip\noindent\textbf{Part II --- Bisection.}
\begin{enumerate}[resume]
\item Check that the invariant $lo^2 \leq n < hi^2$ holds initially.
\item Show that it is preserved, and that $lo < mid < hi$ in every iteration.
\item Show that $hi - lo$ is a variant, and deduce termination.
\item Prove the conclusion: at the exit, $lo$ satisfies the \omterm{def:b1:program-correctness:specification}{specification}.
\item Show that the number of iterations is at most $\lceil \log_2 (n + 1) \rceil$, and compare with the table.
\item In C with 64-bit unsigned integers, which expressions of \code{isqrt_bisect} can overflow for $n = 2^{64} - 1$? How
would you avoid it?
\end{enumerate}

\medskip\noindent\textbf{Part III --- Newton's method on integers.}
\begin{enumerate}[resume]
\item Let $s$ be the integer square root of $n \geq 1$. For an integer $r \geq 1$, show that $\lfloor (r + \lfloor n/r
\rfloor)/2 \rfloor = \lfloor (r + n/r)/2 \rfloor$.
\item Show that $(r + n/r)/2 \geq \sqrt{n}$ for $r > 0$, and deduce that the new $r$ is at least $s$: the invariant
$r \geq s$ is preserved.
\item Show that if $r^2 > n$, the new $r$ is strictly smaller than $r$.
\item Deduce that the loop terminates, and that at the exit $r = s$.
\item Why must the case $n = 0$ be treated apart?
\item Why does the loop test $r^2 > n$ rather than $r \neq$ new $r$?
\end{enumerate}

\medskip\noindent\textbf{Part IV --- Tests and counts.}
\begin{enumerate}[resume]
\item Read the table: for which $n$ does bisection reach $\lceil \log_2 (n + 1) \rceil$, and for which does it stay
below?
\item Newton's method needs 37 iterations for $2^{64} - 1$. Explain why the first iterations roughly halve $r$, and
why the last few are fast.
\item Suggest a better starting value than $n$ and say what must be proved again.
\item The \omterm{def:b1:program-correctness:oracle}{property-based test} draws $n$ at random below $2^{64}$. Why does it also check every $n < 2000$ separately?
\item What would the test have caught if \code{isqrt_newton} returned $r$ with $r^2 \geq n$ rather than $r^2 > n$ in its
loop condition? Check by reasoning on $n = 16$.
\item State the named result of the problem, \emph{Newton's integer square root is \omterm{def:b1:program-correctness:correctness}{totally correct}}, with the measured
iteration counts against bisection's.
\end{enumerate}
\end{problem}
